\subsection{Native correctness and source binding}

All 1,112 maintained tests pass with Regina in 210.025 seconds. Seven new
tests cover random multibatch agreement with literal substitution, negative
handles, forward bindings, empty images, repeated child edges, invalid late
batches, shared node/work allowances and false-valued cancellation callbacks.
Large binary powers are replayed with expansion, equality, reduction, inversion,
slicing and ordinary substitution helpers disabled. A 128-step fixture checks
the historical depth/size recurrence and final raw length $2^{128}+2$.
Actual diagram certificates split into consecutive singleton batches pass both
source verifiers, and explicit normalization boundaries create separate blocks.

The final independent audit takes 9.257 seconds and freezes 431 Python and
corpus hashes. It retains 400 previous single-batch DAG fixtures with 1,600
compressed and 800 literal replays. Another 400 random raw blocks contribute
800 old/new compressed comparisons and 1,151 literal batch steps. All agree
as literal words on every retained slot and on the surviving generator set;
135 of the new random blocks exercise legacy dependency-order recovery.

On 79 validated diagrams, both default and optional-batch producer certificates
are identical to baseline \texttt{63c262ac3c22}; there are no gained or lost
positive outcomes. All 920 original source-bound replay checks pass. Splitting
48 distinct actual source certificates into single-entry batches yields
another 192 old/current literal/compressed checks, all passing; all 48 exercise
the new path, with 49 persistent blocks in total. These split certificates
are valid alternative witnesses for the same sources. The maintained producer
does not generally choose that serial encoding.

\subsection{Complete supplied-block measurements}

The pinned kernel run takes 350.994 seconds and retains 240 measured calls
and 48 warm-ups. Four shuffled arms provide old/new and separate A/A pairs;
each fixture has five measured rounds. Timers include fresh binary grammar
construction and complete checked replay. All calls complete. Exact output
lengths are checked outside the timers; retained structural hashes describe
the resulting grammars and are not word-equality proofs when associations
differ. Small literal comparisons and the source-bound audit provide separate
semantic evidence. Table entries are median milliseconds and median paired
old/new ratios, which need not equal ratios of the separate medians.

\begin{center}\small
\begin{tabular}{@{}lrrrrrr@{}}
\toprule
Family & pivots & old ms & new ms & old/new & old A/A & new A/A \\
\midrule
contexts & 8 & 2.715 & 1.212 & 2.301 & 0.988 & 1.009 \\
contexts & 32 & 164.627 & 27.204 & 6.061 & 1.006 & 0.983 \\
contexts & 128 & 12125.140 & 713.619 & 16.988 & 1.008 & 0.993 \\
tower & 8 & 21.830 & 14.760 & 1.474 & 0.994 & 0.996 \\
tower & 32 & 391.609 & 106.792 & 3.722 & 1.010 & 1.012 \\
tower & 128 & 6597.190 & 1095.113 & 6.156 & 1.001 & 0.994 \\
legacy & 8 & 24.184 & 16.445 & 1.485 & 0.951 & 0.995 \\
legacy & 32 & 335.460 & 127.707 & 2.616 & 0.989 & 0.995 \\
legacy & 128 & 5342.976 & 1541.778 & 3.446 & 1.003 & 0.961 \\
shallow & 8 & 0.351 & 0.331 & 1.083 & 0.985 & 0.985 \\
shallow & 32 & 1.700 & 1.635 & 1.111 & 0.970 & 0.822 \\
shallow & 128 & 3.935 & 3.574 & 1.109 & 0.991 & 0.998 \\
\bottomrule
\end{tabular}
\end{center}


The context family is the delivery's explicit artificial presentation. With
$a=1$, $b=2$, $x_i=i+3$ and $z=k+4$, its relators are
\[
 r_0=x_0^{-1}a x_1b,\qquad
 r_i=x_0^{-1}x_{i+1}\ (1\le i<k),\qquad
 r_k=x_0^{-1}z.
\]
Eliminate $x_i$ from slot $i$ for $i=0,\ldots,k-1$. The retained raw word
has length $2^k+2$. Its padding includes cancelling contexts that an intervening
free reduction would remove; this is not a hard-knot construction.
At $k=128$, complete fresh construction/replay takes 12,125.140 ms under the
old checker and 713.619 ms persistently, with paired ratio 16.988. Replay work
falls from 17,679,028 to 1,366,305. The old checker adds 57,155 ordinary arena
nodes; the new checker adds 12,031 at export and retains 6,890 private nodes
during replay. Its peak combined node count is 19,184, including the 263-node
source arena. The maximum binding-free depth is 855 and the final length has
129 bits. Both engines perform the same 128 raw singleton transformations.

The tower families use $2^{128}$-power contexts and 128 pivots arranged into
32 batches of four. Ordered replay takes 6,597.190 versus 1,095.113 ms,
with paired ratio 6.156 and work 9,094,171 versus 1,888,530. Ordinary nodes
added are nearly equal, 33,437 versus 33,406: avoiding repeated intermediate
traversals is the main saving here. The private graph has 17,026 nodes and
the peak combined count is 67,587, more than the old final ordinary count
50,592. The largest output length needs 16,514 bits. Reversing each batch's
entry order exercises legacy order recovery 32 times; its paired ratio is
3.446 and work falls from 10,971,306 to 2,978,866. These examples show both
the benefit and the additional transient storage of persistent replay.

Two shallow independent batches provide a conversion control. At 128 pivots,
the recorded paired ratio is 1.109, but charged replay work rises from 4,605
to 5,057 and the peak combined node count is 781, versus 394 old ordinary
nodes. At 32 pivots the new-arm A/A ratio is 0.822, so its small timing gain
is especially noisy. Across all kernel fixtures the median A/A ranges are
0.951--1.010 historically and 0.822--1.012 persistently. The result is a
whole-block representation improvement with favorable large-fixture timings,
not a promise of fewer operations or less memory on every input.

\subsection{Supplied source proofs and unchanged discovery}

The separate source benchmark retains 80 measured calls and 16 warm-ups in
23.549 seconds. Each old/new pair verifies the identical version-eight proof
for a stabilized circle diagram, split into consecutive single-entry batches.
Every timed call performs fresh diagram validation and the verifier's complete
independent presentation reconstruction and replay. Producing and splitting
the supplied certificate, serialization and the separate literal cross-check
are outside the timers. All calls complete. These are genuine source-bound
proofs, but their serial encoding is supplied to the verifier rather than
chosen by the current discovery algorithm.

\begin{center}\small
\begin{tabular}{@{}lrrrrr@{}}
\toprule
Crossings & old ms & new ms & old/new & old A/A & new A/A \\
\midrule
8 & 1.360 & 0.868 & 1.573 & 1.069 & 0.994 \\
32 & 16.699 & 7.723 & 2.155 & 0.999 & 0.971 \\
128 & 262.194 & 112.566 & 2.336 & 1.020 & 1.000 \\
256 & 1064.499 & 449.934 & 2.345 & 0.996 & 0.958 \\
\bottomrule
\end{tabular}
\end{center}


At 256 crossings, the unchanged 23,189-byte proof has 255 consecutive raw
batches. Complete verification takes 1,064.499 versus 449.934 ms, with paired
ratio 2.345. The replay arena's work counter falls from 1,637,855 to 894,796;
these counters exclude the preceding independent presentation reconstruction,
which is included in the timed calls. Final ordinary nodes fall from 2,800
to 1,785, while 1,273 private circuit nodes make the new combined peak 3,058.
Thus the time saving does not imply a smaller peak node count. Across the
four fixtures the paired ratios range from 1.573 to 2.345; median A/A ranges
are 0.996--1.069 historically and 0.958--1.000 persistently.

The ordinary group-stage benchmark separately runs fresh source validation,
optional-batch search and mandatory replay. All 120 measured calls and 24
warm-ups finish in 2.436 seconds. Its producer emits a single batch on these
circle fixtures, so the persistent block path does not activate. Paired
ratios range from 0.988 to 1.006; historical A/A ratios range from 0.986 to
1.035 and new A/A ratios from 0.992 to 1.026. This gives no useful discovery
speedup and confirms why the supplied split-proof result must be labelled
separately.

\begin{center}\small
\begin{tabular}{@{}lrrrrr@{}}
\toprule
Crossings & old ms & new ms & old/new & old A/A & new A/A \\
\midrule
8 & 1.258 & 1.250 & 1.006 & 0.986 & 0.992 \\
16 & 2.492 & 2.498 & 1.004 & 1.001 & 1.005 \\
32 & 4.928 & 4.959 & 0.988 & 1.008 & 1.001 \\
64 & 10.011 & 10.029 & 0.994 & 0.995 & 0.999 \\
128 & 20.254 & 20.442 & 0.994 & 1.007 & 1.000 \\
256 & 44.013 & 42.488 & 0.993 & 1.035 & 1.026 \\
\bottomrule
\end{tabular}
\end{center}


The whole-recognition comparison uses the same optional group portfolio in
both engines on nineteen native diagrams. All 380 measured calls and 76
warm-ups complete in 39.931 seconds, retaining the same certificates across
arms. One successful proof, \texttt{survivor-10}, has consecutive batches of
seven and two generators before normalization. Its paired ratio is 1.047,
with old/new medians 5.097 and 4.906 ms; its old/new A/A ratios are 1.007
and 1.002. This is a modest observed gain on an activating native input.
Separate deterministic source replay reduces arena work from 1,904 to 1,783
and final ordinary nodes from 165 to 123. The new block has 57 private nodes
and peaks at 154 combined nodes before the later normalization. These counters
exclude source reconstruction and are retained with an additional literal
check in \path{data/persistent-replay-activation.json}.
The Gordian trial still exhausts its producer allowance and
uses the unchanged discovery fallback; its final proof cannot benefit from
the new block checker. Its paired ratio is 1.001, with separate medians
1,254.964 and 1,273.125 ms. Across the corpus ratios range from 0.979 to
1.047, with old A/A medians 0.987--1.026 and new A/A 0.981--1.015. This
establishes no broad whole-recognition improvement.

\begin{center}\small
\begin{tabular}{@{}lrrrrr@{}}
\toprule
Input & old ms & new ms & old/new & old A/A & new A/A \\
\midrule
survivor-00 & 14.047 & 14.072 & 1.003 & 1.009 & 0.997 \\
survivor-01 & 28.759 & 27.750 & 1.034 & 1.026 & 0.997 \\
survivor-02 & 9.234 & 9.132 & 1.007 & 1.009 & 0.988 \\
survivor-03 & 37.838 & 38.481 & 0.995 & 1.004 & 1.002 \\
mirror-03 & 5.672 & 5.745 & 0.987 & 0.999 & 1.007 \\
survivor-04 & 28.870 & 29.141 & 0.995 & 0.996 & 0.997 \\
survivor-05 & 15.128 & 15.316 & 0.988 & 0.998 & 1.002 \\
survivor-06 & 5.085 & 5.145 & 0.994 & 0.998 & 0.991 \\
survivor-07 & 19.294 & 19.317 & 0.998 & 0.987 & 1.005 \\
survivor-08 & 6.857 & 6.897 & 1.001 & 1.003 & 1.002 \\
mirror-08 & 15.080 & 15.078 & 1.003 & 0.999 & 1.001 \\
survivor-09 & 5.402 & 5.420 & 1.003 & 1.005 & 0.993 \\
survivor-10 & 5.097 & 4.906 & 1.047 & 1.007 & 1.002 \\
survivor-11 & 4.681 & 4.668 & 1.000 & 1.008 & 1.000 \\
gordian & 1254.964 & 1273.125 & 1.001 & 1.004 & 0.987 \\
trefoil & 0.057 & 0.057 & 0.979 & 0.995 & 0.981 \\
figure\_eight & 0.062 & 0.062 & 0.994 & 0.990 & 1.006 \\
conway & 0.978 & 0.991 & 0.987 & 0.994 & 1.015 \\
kinoshita\_terasaka & 1.003 & 1.013 & 0.996 & 0.996 & 0.988 \\
\bottomrule
\end{tabular}
\end{center}


In total the four new benchmark modes retain 820 measured calls and 164
warm-ups, all complete. Kernel, supplied-source, group-discovery and full
recognition scopes are kept separate. None of the runs overlaps tests,
article builds or another benchmark. The driver
\path{../fast/compressed_word_research/persistent_replay.py} provides
\texttt{audit}, \texttt{kernels}, \texttt{stages} and \texttt{pipeline};
\path{data/persistent_replay_source.py} runs the supplied-source comparison.
The former freezes 431 source hashes and the latter additionally freezes
its own script, for 432. The baseline package is isolated from commit
\texttt{63c262ac3c22}; only its compressed group-replay dispatcher differs
among existing runtime modules, plus the new persistent checker. All raw
records, proof objects and reproduction metadata are retained under
\path{data/persistent-replay-*}.
The three deterministic activation checks are reproduced by
\path{data/persistent_replay_activation.py}; its 433 hashes additionally bind
the script and the selected pipeline record.
